\section{Movimiento de proyectiles}

Lanzar una pelota hacia una canasta, patear un balón, dejar caer un paquete desde
un vehículo en movimiento o hacer salir agua de una fuente son situaciones que,
aunque parecen diferentes, comparten una característica importante: una vez que
el objeto queda libre, puede describir una trayectoria curva mientras la gravedad
modifica continuamente su movimiento vertical.

En el capítulo anterior aprendimos a describir movimientos en dos y tres
dimensiones mediante componentes rectangulares. Ahora aplicaremos esas ideas a
uno de los modelos más importantes de la cinemática: el movimiento de
proyectiles.

\begin{idea}[Dos movimientos simultáneos]
La clave para comprender el movimiento de un proyectil consiste en reconocer que
puede analizarse como la combinación de dos movimientos que ocurren
simultáneamente:

\[
\text{movimiento horizontal}
\]

y

\[
\text{movimiento vertical}.
\]

En el modelo ideal:

\[
a_x=0,
\qquad
a_y=-g.
\]

Por tanto, horizontalmente existe movimiento uniforme, mientras que
verticalmente existe movimiento uniformemente acelerado debido a la gravedad.

Ambas componentes están unidas por una variable común:

\[
\boxed{t}.
\]
\end{idea}


\section{¿Qué es un proyectil?}

\begin{definicion}[Proyectil]
Un \textbf{proyectil} es un objeto que, después de recibir una velocidad inicial,
se mueve durante cierto intervalo bajo la acción de la aceleración gravitatoria,
dentro de un modelo en el que se desprecia la resistencia del aire y cualquier
otra interacción significativa.
\end{definicion}

Pueden modelarse como proyectiles, durante parte de su movimiento:

\begin{itemize}
    \item una pelota después de abandonar la mano de quien la lanza;
    \item un balón después de ser pateado;
    \item una piedra después de ser arrojada;
    \item un paquete después de ser liberado desde un vehículo;
    \item una gota de agua lanzada por una fuente;
    \item un objeto que sale horizontalmente del borde de una mesa.
\end{itemize}

La expresión ``después de recibir una velocidad inicial'' es importante.

Mientras una persona impulsa una pelota con la mano, o mientras una máquina
todavía está acelerando un objeto dentro de un mecanismo de lanzamiento, el
objeto todavía no se encuentra en la etapa que modelaremos mediante las
ecuaciones de proyectiles.

\begin{advertencia}
Las ecuaciones del movimiento de proyectiles se aplican \textbf{después de que
el objeto ha sido liberado} y mientras se cumplan las hipótesis del modelo.

No deben utilizarse para describir la fase durante la cual una mano, una rampa,
un resorte u otro mecanismo todavía está impulsando el objeto.
\end{advertencia}


\section{Modelo ideal de proyectil}

El movimiento real de una pelota o de cualquier otro objeto puede estar afectado
por numerosos factores.

Para construir primero una descripción comprensible utilizaremos un modelo
ideal.

\begin{modelo}[Proyectil ideal]
Representaremos el proyectil como una partícula y supondremos que:

\begin{itemize}
    \item la resistencia del aire es despreciable;
    \item no existe viento;
    \item la aceleración gravitatoria es constante;
    \item el movimiento ocurre cerca de la superficie terrestre;
    \item las distancias son suficientemente pequeñas para despreciar la curvatura
    terrestre;
    \item se desprecia la rotación de la Tierra;
    \item después del lanzamiento, la aceleración del proyectil es únicamente la
    aceleración gravitatoria.
\end{itemize}
\end{modelo}

Dentro de este modelo:

\[
\vect{a}=-g\jhat
\]

si elegimos el eje \(+y\) hacia arriba.

Por tanto:

\[
\boxed{a_x=0},
\]

y:

\[
\boxed{a_y=-g}.
\]

Utilizaremos:

\[
g\approx\SI{9.81}{\meter\per\second\squared}.
\]

\begin{interpretacion}
La gravedad no ``empuja hacia adelante'' ni ``frena horizontalmente'' al
proyectil dentro de este modelo.

Su aceleración es vertical.

Por ello:

\[
a_x=0
\]

aunque el proyectil continúe avanzando horizontalmente.
\end{interpretacion}


\section{Sistema de referencia}

Adoptaremos normalmente:

\begin{itemize}
    \item el eje \(x\) horizontal;
    \item el eje \(y\) vertical;
    \item \(+x\) en el sentido horizontal inicial del movimiento;
    \item \(+y\) hacia arriba.
\end{itemize}

Con esta elección:

\[
\vect{a}
=
-g\jhat.
\]

\begin{center}
\begin{tikzpicture}[scale=0.95]

    \draw[-{Latex},thick] (-0.5,0) -- (7.2,0)
        node[right] {\(x\)};

    \draw[-{Latex},thick] (0,-0.5) -- (0,5.2)
        node[above] {\(y\)};

    \draw[
        very thick,
        AzulPrincipal,
        domain=0:6.2,
        samples=80
    ]
    plot (\x,{0.95*\x-0.13*\x*\x});

    \fill (0,0) circle (2.5pt);
    \node[below left] at (0,0) {\(O\)};

    \draw[-{Latex},very thick,VerdePrincipal]
        (0,0) -- (2.1,2.4)
        node[above right] {\(\vect{v}_0\)};

    \draw[dashed] (0,0) -- (2.1,0);
    \draw[dashed] (2.1,0) -- (2.1,2.4);

    \draw[-{Latex},thick,AzulSecundario]
        (0,0.25) -- (1.8,0.25)
        node[midway,below] {\(v_{0x}\)};

    \draw[-{Latex},thick,AzulSecundario]
        (0.25,0) -- (0.25,2.0)
        node[midway,left] {\(v_{0y}\)};

    \draw[-{Latex},very thick,RojoPrincipal]
        (4.8,3.2) -- (4.8,2.0)
        node[right] {\(\vect{a}=-g\jhat\)};

    \draw (0.75,0) arc[start angle=0,end angle=48,radius=0.75];
    \node at (1.0,0.37) {\(\theta_0\)};

\end{tikzpicture}
\end{center}

\begin{advertencia}
El signo:

\[
a_y=-g
\]

no significa que \(g\) sea una cantidad negativa.

La magnitud:

\[
g>0
\]

y el signo negativo aparece porque hemos definido el sentido \(+y\) hacia
arriba, mientras la aceleración gravitatoria apunta hacia abajo.
\end{advertencia}


\section{Velocidad inicial y componentes rectangulares}

Supongamos que el proyectil se lanza con rapidez inicial:

\[
v_0
\]

formando un ángulo:

\[
\theta_0
\]

respecto del eje horizontal positivo.

El vector velocidad inicial es:

\[
\vect{v}_0
=
v_{0x}\ihat+v_{0y}\jhat.
\]

Utilizando componentes rectangulares:

\[
\boxed{
v_{0x}=v_0\cos\theta_0
},
\]

\[
\boxed{
v_{0y}=v_0\sen\theta_0
}.
\]

Por tanto:

\[
\boxed{
\vect{v}_0
=
v_0\cos\theta_0\,\ihat
+
v_0\sen\theta_0\,\jhat
}.
\]

\begin{recordatorio}[Método de componentes]
En este proyecto utilizaremos prioritariamente componentes rectangulares.

El procedimiento será:

\[
\vect{v}_0
\longrightarrow
v_{0x},v_{0y}
\longrightarrow
\text{resolver cada eje}
\longrightarrow
\text{reconstruir el resultado}.
\]

Este método es general y evita depender de fórmulas particulares para cada
configuración geométrica.
\end{recordatorio}


\section{La independencia de las componentes}

Una vez conocidas las componentes iniciales, analizamos por separado los ejes
\(x\) y \(y\).

En la dirección horizontal:

\[
a_x=0.
\]

En la dirección vertical:

\[
a_y=-g.
\]

Por ello:

\[
\boxed{
\text{horizontal}
\longrightarrow
\text{MRU}
}
\]

mientras:

\[
\boxed{
\text{vertical}
\longrightarrow
\text{MRUA}
}
\]

con aceleración:

\[
-g.
\]

\begin{principio}[Superposición de movimientos]
Dentro del modelo ideal de proyectil, el movimiento bidimensional puede
analizarse como la superposición de:

\begin{itemize}
    \item un movimiento horizontal con velocidad constante;
    \item un movimiento vertical de caída libre.
\end{itemize}

Ambas componentes evolucionan simultáneamente y comparten el mismo tiempo.
\end{principio}

\begin{interpretacion}
El movimiento horizontal no hace que el objeto ``tarde más en caer''.

Si dos objetos comienzan desde la misma altura y con la misma componente
vertical inicial, uno soltado directamente y otro con velocidad horizontal,
ambos tienen exactamente el mismo movimiento vertical dentro del modelo ideal.

La componente horizontal sólo determina cuánto avanzará uno de ellos mientras
cae.
\end{interpretacion}


\section{Ecuaciones horizontales}

Como:

\[
a_x=0,
\]

la componente horizontal de la velocidad permanece constante:

\[
v_x=v_{0x}.
\]

Por tanto:

\[
\boxed{
v_x=v_0\cos\theta_0
}.
\]

La ecuación de posición horizontal es:

\[
x=x_0+v_{0x}t.
\]

Así:

\[
\boxed{
x=x_0+(v_0\cos\theta_0)t
}.
\]

El desplazamiento horizontal es:

\[
\boxed{
\Delta x=v_{0x}t
}.
\]

\begin{interpretacion}
En intervalos iguales de tiempo, el proyectil recorre desplazamientos
horizontales iguales.

Por ejemplo, si:

\[
v_x=\SI{10}{\meter\per\second},
\]

entonces durante cada segundo avanza horizontalmente:

\[
\SI{10}{\meter},
\]

independientemente de si en ese instante está subiendo o bajando.
\end{interpretacion}


\section{Ecuaciones verticales}

En la dirección vertical:

\[
a_y=-g.
\]

La velocidad vertical es:

\[
v_y=v_{0y}-gt.
\]

Por tanto:

\[
\boxed{
v_y=v_0\sen\theta_0-gt
}.
\]

La posición vertical es:

\[
y=y_0+v_{0y}t-\frac{1}{2}gt^2.
\]

Entonces:

\[
\boxed{
y
=
y_0
+
(v_0\sen\theta_0)t
-
\frac{1}{2}gt^2
}.
\]

También podemos utilizar:

\[
v_y^2
=
v_{0y}^2
-
2g(y-y_0).
\]

Es decir:

\[
\boxed{
v_y^2
=
(v_0\sen\theta_0)^2
-
2g(y-y_0)
}.
\]

Estas son exactamente las ecuaciones de caída libre desarrolladas
anteriormente.


\section{Ecuaciones generales del movimiento}

Reuniendo ambos ejes:

\[
\boxed{
x
=
x_0
+
v_0\cos\theta_0\,t
},
\]

\[
\boxed{
y
=
y_0
+
v_0\sen\theta_0\,t
-
\frac{1}{2}gt^2
}.
\]

Las componentes de velocidad son:

\[
\boxed{
v_x=v_0\cos\theta_0
},
\]

\[
\boxed{
v_y=v_0\sen\theta_0-gt
}.
\]

El vector velocidad es:

\[
\boxed{
\vect{v}
=
v_0\cos\theta_0\,\ihat
+
\left(
v_0\sen\theta_0-gt
\right)\jhat
}.
\]

El vector posición es:

\[
\boxed{
\vect{r}
=
\left[
x_0+v_0\cos\theta_0\,t
\right]\ihat
+
\left[
y_0+v_0\sen\theta_0\,t-\frac{1}{2}gt^2
\right]\jhat
}.
\]

Y la aceleración:

\[
\boxed{
\vect{a}=-g\jhat
}.
\]


\section{Rapidez y dirección de la velocidad}

Aunque:

\[
v_x
\]

permanece constante, la componente:

\[
v_y
\]

cambia continuamente.

La rapidez instantánea es:

\[
v
=
\norm{\vect{v}}.
\]

Entonces:

\[
\boxed{
v
=
\sqrt{v_x^2+v_y^2}
}.
\]

Sustituyendo las componentes:

\[
\boxed{
v
=
\sqrt{
(v_0\cos\theta_0)^2
+
(v_0\sen\theta_0-gt)^2
}
}.
\]

La dirección instantánea de la velocidad puede obtenerse mediante:

\[
\tan\theta
=
\frac{v_y}{v_x}.
\]

Por tanto:

\[
\boxed{
\theta
=
\arctan
\left(
\frac{v_y}{v_x}
\right)
}
\]

teniendo siempre en cuenta el cuadrante físico correcto.

\begin{interpretacion}
El vector velocidad es tangente a la trayectoria.

Durante el ascenso:

\[
v_y>0.
\]

En el punto más alto:

\[
v_y=0.
\]

Durante el descenso:

\[
v_y<0.
\]

Sin embargo, mientras el proyectil continúe avanzando hacia \(+x\):

\[
v_x>0
\]

durante todo el vuelo.
\end{interpretacion}


\section{El punto más alto}

Cuando un proyectil se lanza con componente vertical inicial positiva, asciende
hasta alcanzar un punto máximo.

En ese instante:

\[
\boxed{v_y=0}.
\]

Sin embargo:

\[
v_x\neq0
\]

en general.

Por tanto, el proyectil no está completamente en reposo.

\begin{errorfrecuente}[Velocidad en la altura máxima]
En el punto más alto no debe escribirse:

\[
\vect{v}=\vect{0}.
\]

Lo correcto es:

\[
v_y=0,
\]

mientras:

\[
v_x=v_0\cos\theta_0.
\]

Por tanto:

\[
\vect{v}
=
v_x\ihat.
\]

La velocidad es horizontal en ese instante.
\end{errorfrecuente}

La aceleración tampoco desaparece:

\[
\boxed{
\vect{a}=-g\jhat
}.
\]

Así, en el punto más alto:

\[
\vect{v}\perp\vect{a}.
\]


\section{Tiempo para alcanzar la altura máxima}

Partimos de:

\[
v_y=v_{0y}-gt.
\]

En la altura máxima:

\[
v_y=0.
\]

Entonces:

\[
0=v_{0y}-gt_H.
\]

Por tanto:

\[
\boxed{
t_H=\frac{v_{0y}}{g}
}.
\]

Como:

\[
v_{0y}=v_0\sen\theta_0,
\]

obtenemos:

\[
\boxed{
t_H
=
\frac{v_0\sen\theta_0}{g}
}.
\]

Esta ecuación supone:

\[
v_{0y}>0.
\]


\section{Altura máxima}

Utilizamos la relación vertical:

\[
v_y^2
=
v_{0y}^2
-
2g(y-y_0).
\]

En el punto más alto:

\[
v_y=0.
\]

Entonces:

\[
0
=
v_{0y}^2
-
2g(y_H-y_0).
\]

Por tanto:

\[
y_H-y_0
=
\frac{v_{0y}^2}{2g}.
\]

Definiendo:

\[
H=y_H-y_0,
\]

como la altura ganada respecto del punto de lanzamiento:

\[
\boxed{
H
=
\frac{v_{0y}^2}{2g}
}.
\]

En términos de \(v_0\) y \(\theta_0\):

\[
\boxed{
H
=
\frac{
v_0^2\sen^2\theta_0
}{
2g
}
}.
\]

\begin{advertencia}
Esta \(H\) representa la altura máxima \textbf{respecto del punto de lanzamiento}.

Si el proyectil se lanza desde una altura inicial:

\[
y_0\neq0,
\]

la coordenada máxima es:

\[
\boxed{
y_H
=
y_0
+
\frac{
v_0^2\sen^2\theta_0
}{
2g
}
}.
\]
\end{advertencia}


\section{La trayectoria es parabólica}

Hasta ahora hemos descrito:

\[
x(t)
\]

y:

\[
y(t).
\]

Podemos eliminar el tiempo y obtener directamente la relación entre \(x\) y
\(y\).

Partimos de:

\[
x
=
x_0
+
v_0\cos\theta_0\,t.
\]

Despejando:

\[
t
=
\frac{x-x_0}
{v_0\cos\theta_0}.
\]

Sustituimos en:

\[
y
=
y_0
+
v_0\sen\theta_0\,t
-
\frac{1}{2}gt^2.
\]

Entonces:

\[
y
=
y_0
+
v_0\sen\theta_0
\left(
\frac{x-x_0}
{v_0\cos\theta_0}
\right)
-
\frac{1}{2}g
\left(
\frac{x-x_0}
{v_0\cos\theta_0}
\right)^2.
\]

Simplificando:

\[
\boxed{
y
=
y_0
+
(x-x_0)\tan\theta_0
-
\frac{
g(x-x_0)^2
}{
2v_0^2\cos^2\theta_0
}
}.
\]

Esta expresión tiene la forma:

\[
y=A+Bx-Cx^2,
\]

que corresponde a una parábola.

\begin{principio}[Trayectoria ideal]
Dentro del modelo de proyectil ideal, con aceleración gravitatoria uniforme y
sin resistencia del aire, la trayectoria en un plano vertical es parabólica.
\end{principio}

\begin{advertencia}
Una pelota real no sigue necesariamente una parábola perfecta.

Cuando la resistencia del aire, el viento o la rotación del objeto tienen
efectos importantes, la trayectoria puede diferir considerablemente de la
predicción ideal.
\end{advertencia}


\section{Diagrama de movimiento}

Una manera útil de comprender el movimiento consiste en representar la posición
del proyectil en intervalos iguales de tiempo.

\begin{center}
\begin{tikzpicture}[scale=0.95]

    \draw[-{Latex},thick] (-0.4,0) -- (7.2,0)
        node[right] {\(x\)};

    \draw[-{Latex},thick] (0,-0.3) -- (0,4.4)
        node[above] {\(y\)};

    \draw[
        very thick,
        AzulPrincipal,
        domain=0:6.2,
        samples=100
    ]
    plot (\x,{1.05*\x-0.17*\x*\x});

    \foreach \x in {0,1,2,3,4,5,6}{
        \pgfmathsetmacro{\yy}{1.05*\x-0.17*\x*\x}
        \fill (\x,\yy) circle (2.2pt);
    }

    \draw[-{Latex},thick,RojoPrincipal]
        (1,1.55) -- (1,0.80);

    \draw[-{Latex},thick,RojoPrincipal]
        (3,2.55) -- (3,1.80);

    \draw[-{Latex},thick,RojoPrincipal]
        (5,1.55) -- (5,0.80);

    \node[right] at (5,1.15) {\(\vect{a}\)};

\end{tikzpicture}
\end{center}

Observe que los puntos están separados horizontalmente por cantidades iguales si
los intervalos de tiempo son iguales.

Esto ocurre porque:

\[
v_x=\text{constante}.
\]

En cambio, las separaciones verticales cambian debido a:

\[
a_y=-g.
\]


\section{Lanzamiento horizontal}

Un caso importante ocurre cuando el proyectil abandona una superficie con
velocidad completamente horizontal.

Entonces:

\[
\theta_0=0.
\]

Por tanto:

\[
v_{0x}=v_0,
\]

y:

\[
v_{0y}=0.
\]

Las ecuaciones se reducen a:

\[
\boxed{
x=x_0+v_0t
},
\]

\[
\boxed{
y=y_0-\frac{1}{2}gt^2
},
\]

\[
\boxed{
v_x=v_0
},
\]

\[
\boxed{
v_y=-gt
}.
\]

\begin{idea}[El objeto comienza a caer inmediatamente]
Un proyectil lanzado horizontalmente no permanece horizontal durante cierto
tiempo antes de comenzar a caer.

Desde el primer instante después de ser liberado:

\[
a_y=-g.
\]

Por tanto, su componente vertical de velocidad comienza a cambiar
inmediatamente.
\end{idea}


\section{Lanzamiento horizontal desde una altura}

Supongamos que un objeto se lanza horizontalmente desde una altura \(h\) sobre
el suelo.

Elegimos:

\[
y_0=h,
\qquad
y_f=0.
\]

Como:

\[
v_{0y}=0,
\]

tenemos:

\[
0
=
h-\frac{1}{2}gt^2.
\]

Por tanto:

\[
\boxed{
t_f
=
\sqrt{\frac{2h}{g}}
}.
\]

Observe que este tiempo no depende de la velocidad horizontal inicial.

El desplazamiento horizontal hasta el impacto es:

\[
R=v_0t_f.
\]

Entonces:

\[
\boxed{
R
=
v_0
\sqrt{\frac{2h}{g}}
}.
\]

\begin{interpretacion}
Para una altura fija, aumentar la rapidez horizontal no modifica el tiempo que
tarda el proyectil en llegar al suelo dentro del modelo ideal.

Lo que aumenta es la distancia horizontal recorrida durante ese mismo tiempo.
\end{interpretacion}


\section{Ejemplo: lanzamiento horizontal}

\begin{ejemplo}[Objeto que sale del borde de una plataforma]
Una pelota abandona horizontalmente una plataforma situada a:

\[
h=\SI{20.0}{\meter}
\]

sobre el suelo con rapidez:

\[
v_0=\SI{12.0}{\meter\per\second}.
\]

Despreciando la resistencia del aire, determine:

\begin{enumerate}[label=\alph*)]
    \item el tiempo de caída;
    \item la distancia horizontal recorrida;
    \item las componentes de la velocidad inmediatamente antes del impacto;
    \item la rapidez de impacto;
    \item la dirección de la velocidad al impactar.
\end{enumerate}
\end{ejemplo}

\subsubsection*{Sistema de referencia}

Elegimos:

\[
+x
\]

horizontal en el sentido del lanzamiento y:

\[
+y
\]

hacia arriba.

Tomamos:

\[
x_0=0,
\qquad
y_0=\SI{20.0}{\meter}.
\]

La velocidad inicial tiene componentes:

\[
v_{0x}=\SI{12.0}{\meter\per\second},
\]

\[
v_{0y}=0.
\]


\subsubsection*{Tiempo de caída}

Verticalmente:

\[
y
=
y_0-\frac{1}{2}gt^2.
\]

Al llegar al suelo:

\[
y=0.
\]

Entonces:

\[
0
=
\SI{20.0}{\meter}
-
\frac{1}{2}
(\SI{9.81}{\meter\per\second\squared})t^2.
\]

Por tanto:

\[
t
=
\sqrt{
\frac{
2(\SI{20.0}{\meter})
}{
\SI{9.81}{\meter\per\second\squared}
}
}.
\]

Así:

\[
\boxed{
t\approx\SI{2.02}{\second}
}.
\]


\subsubsection*{Desplazamiento horizontal}

Horizontalmente:

\[
x=v_{0x}t.
\]

Entonces:

\[
x
=
(\SI{12.0}{\meter\per\second})
(\SI{2.02}{\second}).
\]

Por tanto:

\[
\boxed{
x\approx\SI{24.2}{\meter}
}.
\]


\subsubsection*{Velocidad de impacto}

La componente horizontal permanece constante:

\[
\boxed{
v_x=\SI{12.0}{\meter\per\second}
}.
\]

La componente vertical es:

\[
v_y=-gt.
\]

Entonces:

\[
v_y
=
-
(\SI{9.81}{\meter\per\second\squared})
(\SI{2.02}{\second}).
\]

Así:

\[
\boxed{
v_y\approx-\SI{19.8}{\meter\per\second}
}.
\]

El signo negativo indica que la componente vertical apunta hacia abajo.

La rapidez es:

\[
v
=
\sqrt{v_x^2+v_y^2}.
\]

Por tanto:

\[
v
=
\sqrt{
(12.0)^2+(19.8)^2
}
\,
\si{\meter\per\second}.
\]

Así:

\[
\boxed{
v\approx\SI{23.2}{\meter\per\second}
}.
\]


\subsubsection*{Dirección de impacto}

La magnitud del ángulo debajo de la horizontal satisface:

\[
\tan\theta
=
\frac{|v_y|}{v_x}.
\]

Entonces:

\[
\theta
=
\arctan
\left(
\frac{19.8}{12.0}
\right).
\]

Por tanto:

\[
\boxed{
\theta\approx58.8\grado
}
\]

por debajo del eje \(+x\).

\begin{interpretacion}
La pelota mantuvo durante todo el vuelo su componente horizontal de velocidad,
pero adquirió una componente vertical cada vez mayor en magnitud.

Por ello, el vector velocidad fue cambiando de dirección continuamente.
\end{interpretacion}


\section{Lanzamiento oblicuo}

Cuando:

\[
0<\theta_0<90\grado,
\]

la velocidad inicial posee simultáneamente una componente horizontal y una
componente vertical positiva:

\[
v_{0x}=v_0\cos\theta_0,
\]

\[
v_{0y}=v_0\sen\theta_0.
\]

El proyectil comienza ascendiendo porque:

\[
v_{0y}>0.
\]

Durante el ascenso, la gravedad reduce progresivamente \(v_y\).

En la altura máxima:

\[
v_y=0.
\]

Después:

\[
v_y<0,
\]

y el proyectil desciende.

Mientras tanto:

\[
v_x=v_{0x}
\]

permanece constante.


\section{Caso simétrico: mismo nivel de lanzamiento y llegada}

Consideremos ahora una situación especial y muy útil.

El proyectil parte desde:

\[
y_0
\]

y finalmente regresa a esa misma altura:

\[
y_f=y_0.
\]

Además:

\[
v_{0y}>0.
\]

Dentro del modelo ideal aparece una simetría importante.


\subsection{Simetría temporal}

El tiempo empleado en subir desde la altura de lanzamiento hasta la altura
máxima es igual al tiempo empleado en regresar desde la altura máxima hasta esa
misma altura.

Por tanto:

\[
t_{\mathrm{vuelo}}
=
2t_H.
\]

Como:

\[
t_H
=
\frac{v_0\sen\theta_0}{g},
\]

resulta:

\[
\boxed{
t_{\mathrm{vuelo}}
=
\frac{
2v_0\sen\theta_0
}{
g
}
}.
\]

\begin{condicion}
Esta fórmula requiere:

\[
y_f=y_0.
\]

No es una expresión general para cualquier movimiento de proyectil.
\end{condicion}


\subsection{Simetría de las velocidades}

Al regresar a la misma altura:

\[
v_x=v_{0x},
\]

y:

\[
v_y=-v_{0y}.
\]

Por tanto, la rapidez es:

\[
v
=
\sqrt{
v_{0x}^2+(-v_{0y})^2
}.
\]

Así:

\[
\boxed{
v=v_0
}.
\]

\begin{interpretacion}
En ausencia de resistencia del aire, un proyectil que regresa a su altura inicial
lo hace con la misma rapidez con la que fue lanzado, pero con una componente
vertical de sentido contrario.
\end{interpretacion}


\section{Alcance horizontal}

Llamaremos alcance horizontal \(R\) a la distancia horizontal recorrida cuando el
proyectil vuelve a la misma altura desde la que fue lanzado.

Como horizontalmente:

\[
R=v_{0x}t_{\mathrm{vuelo}},
\]

tenemos:

\[
R
=
v_0\cos\theta_0
\left(
\frac{
2v_0\sen\theta_0
}{
g
}
\right).
\]

Por tanto:

\[
R
=
\frac{
2v_0^2\sen\theta_0\cos\theta_0
}{
g
}.
\]

Utilizando:

\[
2\sen\theta_0\cos\theta_0
=
\sen(2\theta_0),
\]

obtenemos:

\[
\boxed{
R
=
\frac{
v_0^2\sen(2\theta_0)
}{
g
}
}.
\]

\begin{condicion}[Condiciones de la fórmula del alcance]
La expresión:

\[
R
=
\frac{
v_0^2\sen(2\theta_0)
}{
g
}
\]

es válida cuando:

\begin{itemize}
    \item se desprecia la resistencia del aire;
    \item \(g\) puede considerarse constante;
    \item el punto final está a la misma altura que el punto inicial;
    \item la superficie de referencia es horizontal.
\end{itemize}
\end{condicion}

\begin{errorfrecuente}
No debe utilizarse automáticamente la fórmula:

\[
R
=
\frac{
v_0^2\sen(2\theta_0)
}{
g
}
\]

cuando el proyectil se lanza desde una montaña, edificio, plataforma u otra
altura diferente de la posición final.

En esos casos debe utilizarse primero la ecuación vertical general para
determinar el tiempo.
\end{errorfrecuente}


\section{Ángulo de alcance máximo}

Para una rapidez inicial fija:

\[
v_0=\text{constante},
\]

y bajo las condiciones del caso simétrico:

\[
R
=
\frac{v_0^2}{g}
\sen(2\theta_0).
\]

El valor máximo de:

\[
\sen(2\theta_0)
\]

es:

\[
1.
\]

Esto ocurre cuando:

\[
2\theta_0=90\grado.
\]

Por tanto:

\[
\boxed{
\theta_0=45\grado
}.
\]

El alcance máximo es entonces:

\[
\boxed{
R_{\max}
=
\frac{v_0^2}{g}
}.
\]

\begin{advertencia}
La afirmación:

\[
\text{``el alcance máximo ocurre siempre a }45\grado\text{''}
\]

es incorrecta.

El resultado de \(45\grado\) corresponde al modelo ideal cuando lanzamiento y
llegada ocurren a la misma altura y se desprecia la resistencia del aire.

Si las alturas son distintas o el aire tiene efectos importantes, el ángulo que
maximiza el alcance puede ser diferente.
\end{advertencia}


\section{Ángulos complementarios}

Para el caso simétrico:

\[
R
=
\frac{
v_0^2\sen(2\theta_0)
}{
g
}.
\]

Consideremos dos ángulos:

\[
\theta_1
\]

y:

\[
\theta_2=90\grado-\theta_1.
\]

Entonces:

\[
2\theta_2
=
180\grado-2\theta_1.
\]

Como:

\[
\sen(180\grado-\alpha)=\sen\alpha,
\]

ambos producen el mismo alcance.

Por tanto:

\[
\boxed{
R(\theta)
=
R(90\grado-\theta)
}.
\]

Por ejemplo, dentro del modelo ideal:

\[
30\grado
\]

y:

\[
60\grado
\]

producen el mismo alcance horizontal si \(v_0\) es igual y el proyectil regresa a
la altura de lanzamiento.

\begin{interpretacion}
Tener el mismo alcance no significa tener la misma trayectoria.

El lanzamiento con mayor ángulo:

\begin{itemize}
    \item alcanza mayor altura;
    \item permanece más tiempo en el aire.
\end{itemize}

El de menor ángulo:

\begin{itemize}
    \item alcanza menor altura;
    \item posee mayor componente horizontal inicial.
\end{itemize}
\end{interpretacion}


\section{Ejemplo: lanzamiento oblicuo al mismo nivel}

\begin{ejemplo}[Pelota lanzada desde el suelo]
Una pelota se lanza desde el suelo con rapidez inicial:

\[
v_0=\SI{20.0}{\meter\per\second}
\]

formando un ángulo:

\[
\theta_0=35.0\grado
\]

con la horizontal.

Despreciando la resistencia del aire, determine:

\begin{enumerate}[label=\alph*)]
    \item las componentes iniciales de la velocidad;
    \item el tiempo para alcanzar la altura máxima;
    \item la altura máxima;
    \item el tiempo total de vuelo;
    \item el alcance horizontal.
\end{enumerate}
\end{ejemplo}


\subsubsection*{Componentes de la velocidad inicial}

Utilizamos:

\[
v_{0x}=v_0\cos\theta_0.
\]

Entonces:

\[
v_{0x}
=
(\SI{20.0}{\meter\per\second})
\cos(35.0\grado).
\]

Por tanto:

\[
\boxed{
v_{0x}\approx\SI{16.4}{\meter\per\second}
}.
\]

La componente vertical es:

\[
v_{0y}=v_0\sen\theta_0.
\]

Así:

\[
v_{0y}
=
(\SI{20.0}{\meter\per\second})
\sen(35.0\grado).
\]

Entonces:

\[
\boxed{
v_{0y}\approx\SI{11.5}{\meter\per\second}
}.
\]


\subsubsection*{Tiempo hasta la altura máxima}

\[
t_H
=
\frac{v_{0y}}{g}.
\]

Por tanto:

\[
t_H
=
\frac{
\SI{11.5}{\meter\per\second}
}{
\SI{9.81}{\meter\per\second\squared}
}.
\]

Así:

\[
\boxed{
t_H\approx\SI{1.17}{\second}
}.
\]


\subsubsection*{Altura máxima}

\[
H
=
\frac{v_{0y}^2}{2g}.
\]

Entonces:

\[
H
=
\frac{
(\SI{11.5}{\meter\per\second})^2
}{
2(\SI{9.81}{\meter\per\second\squared})
}.
\]

Por tanto:

\[
\boxed{
H\approx\SI{6.71}{\meter}
}.
\]


\subsubsection*{Tiempo total de vuelo}

Como regresa al mismo nivel:

\[
t_{\mathrm{vuelo}}=2t_H.
\]

Entonces:

\[
\boxed{
t_{\mathrm{vuelo}}
\approx
\SI{2.34}{\second}
}.
\]


\subsubsection*{Alcance}

Utilizamos:

\[
R=v_{0x}t_{\mathrm{vuelo}}.
\]

Así:

\[
R
=
(\SI{16.4}{\meter\per\second})
(\SI{2.34}{\second}).
\]

Por tanto:

\[
\boxed{
R\approx\SI{38.3}{\meter}
}.
\]

También podríamos utilizar, porque el proyectil vuelve al mismo nivel:

\[
R
=
\frac{
v_0^2\sen(2\theta_0)
}{
g
}.
\]

\begin{interpretacion}
El movimiento horizontal continuó uniformemente durante los
\(\SI{2.34}{\second}\) de vuelo.

Simultáneamente, la componente vertical disminuyó desde un valor positivo hasta
cero y posteriormente se hizo negativa durante el descenso.
\end{interpretacion}


\section{Cuando las alturas inicial y final son diferentes}

Muchos problemas reales no presentan simetría.

Por ejemplo:

\begin{itemize}
    \item un balón puede lanzarse desde una plataforma y caer al suelo;
    \item una piedra puede lanzarse desde un acantilado;
    \item un objeto puede lanzarse desde el suelo hacia una ventana elevada;
    \item un paquete puede liberarse desde una aeronave.
\end{itemize}

En estas situaciones:

\[
y_f\neq y_0.
\]

La estrategia más segura consiste en utilizar siempre las ecuaciones generales.


\subsection{Determinación del tiempo}

La ecuación vertical es:

\[
y_f
=
y_0
+
v_{0y}t
-
\frac{1}{2}gt^2.
\]

Definiendo:

\[
\Delta y=y_f-y_0,
\]

tenemos:

\[
\Delta y
=
v_{0y}t
-
\frac{1}{2}gt^2.
\]

Reordenando:

\[
\frac{1}{2}gt^2
-
v_{0y}t
+
\Delta y
=
0.
\]

Esta ecuación puede resolverse mediante la fórmula cuadrática.

\begin{interpretacion}
El tiempo suele obtenerse primero a partir del movimiento vertical.

Una vez conocido, ese mismo valor de \(t\) se utiliza en el movimiento horizontal:

\[
x_f=x_0+v_{0x}t.
\]

No existen un ``tiempo horizontal'' y un ``tiempo vertical'' distintos.
\end{interpretacion}


\section{Ejemplo: lanzamiento desde una altura}

\begin{ejemplo}[Proyectil lanzado desde una plataforma]
Desde una plataforma situada a:

\[
y_0=\SI{12.0}{\meter}
\]

sobre el suelo se lanza una pelota con rapidez:

\[
v_0=\SI{18.0}{\meter\per\second}
\]

y ángulo:

\[
\theta_0=30.0\grado.
\]

Determine:

\begin{enumerate}[label=\alph*)]
    \item el tiempo que tarda en llegar al suelo;
    \item la distancia horizontal recorrida;
    \item la velocidad inmediatamente antes del impacto;
    \item la rapidez de impacto.
\end{enumerate}
\end{ejemplo}


\subsubsection*{Componentes iniciales}

Horizontalmente:

\[
v_{0x}
=
v_0\cos\theta_0.
\]

Entonces:

\[
v_{0x}
=
(\SI{18.0}{\meter\per\second})
\cos(30.0\grado).
\]

Así:

\[
\boxed{
v_{0x}\approx\SI{15.6}{\meter\per\second}
}.
\]

Verticalmente:

\[
v_{0y}
=
v_0\sen\theta_0.
\]

Por tanto:

\[
\boxed{
v_{0y}=\SI{9.00}{\meter\per\second}
}.
\]


\subsubsection*{Tiempo hasta el suelo}

Tomamos:

\[
y_f=0.
\]

Entonces:

\[
0
=
12.0
+
9.00t
-
4.905t^2.
\]

Reordenando:

\[
4.905t^2
-
9.00t
-
12.0
=
0.
\]

La fórmula cuadrática produce dos raíces.

La raíz positiva es:

\[
\boxed{
t\approx\SI{2.73}{\second}
}.
\]

La otra raíz es negativa y no pertenece al intervalo posterior al lanzamiento.


\subsubsection*{Distancia horizontal}

\[
\Delta x=v_{0x}t.
\]

Entonces:

\[
\Delta x
=
(\SI{15.6}{\meter\per\second})
(\SI{2.73}{\second}).
\]

Por tanto:

\[
\boxed{
\Delta x\approx\SI{42.6}{\meter}
}.
\]


\subsubsection*{Velocidad de impacto}

La componente horizontal es:

\[
\boxed{
v_x\approx\SI{15.6}{\meter\per\second}
}.
\]

La componente vertical:

\[
v_y=v_{0y}-gt.
\]

Así:

\[
v_y
=
\SI{9.00}{\meter\per\second}
-
(\SI{9.81}{\meter\per\second\squared})
(\SI{2.73}{\second}).
\]

Por tanto:

\[
\boxed{
v_y\approx-\SI{17.8}{\meter\per\second}
}.
\]

El vector velocidad es:

\[
\boxed{
\vect{v}
\approx
(15.6\ihat-17.8\jhat)
\,\si{\meter\per\second}
}.
\]

La rapidez es:

\[
v
=
\sqrt{
(15.6)^2+(17.8)^2
}
\,\si{\meter\per\second}.
\]

Entonces:

\[
\boxed{
v\approx\SI{23.7}{\meter\per\second}
}.
\]

\begin{interpretacion}
La rapidez de impacto es mayor que la rapidez inicial porque el proyectil termina
a una altura menor que aquella desde la que fue lanzado.

Además, este problema no es simétrico:

\[
y_f\neq y_0.
\]

Por ello, no era válido utilizar directamente la fórmula especial del alcance.
\end{interpretacion}


\section{Una misma altura puede alcanzarse dos veces}

Cuando un proyectil asciende y luego desciende, puede atravesar una misma altura
en dos instantes diferentes.

Supongamos:

\[
y_0<y<y_H.
\]

La ecuación:

\[
y
=
y_0
+
v_{0y}t
-
\frac{1}{2}gt^2
\]

puede producir dos valores positivos de \(t\).

Uno corresponde al ascenso y el otro al descenso.

En esos dos puntos:

\[
v_x
\]

es el mismo, pero:

\[
v_y
\]

tiene signos opuestos si la altura está a ambos lados simétricos del ápice.

\begin{interpretacion}
Conocer únicamente la posición del proyectil no determina completamente su
estado de movimiento.

También necesitamos conocer su velocidad o saber si está ascendiendo o
descendiendo.
\end{interpretacion}


\section{Simetría de la trayectoria ideal}

Cuando lanzamiento y llegada ocurren a la misma altura, la trayectoria presenta
simetría respecto de la vertical que pasa por el punto más alto.

Para posiciones situadas a igual altura antes y después del ápice:

\[
v_x=\text{constante},
\]

mientras:

\[
v_{y,\mathrm{descenso}}
=
-v_{y,\mathrm{ascenso}}.
\]

Por tanto:

\[
v_{\mathrm{descenso}}
=
v_{\mathrm{ascenso}}.
\]

También:

\[
x_H=\frac{R}{2}.
\]

\begin{advertencia}
Esta simetría pertenece al modelo ideal y al caso:

\[
y_f=y_0.
\]

La resistencia del aire rompe esta simetría y hace que el movimiento real de
muchos proyectiles sea diferente durante el ascenso y el descenso.
\end{advertencia}


\section{Comparación de distintos ángulos de lanzamiento}

Para una misma rapidez inicial, el ángulo modifica las dos componentes:

\[
v_{0x}=v_0\cos\theta_0,
\]

\[
v_{0y}=v_0\sen\theta_0.
\]

Si aumentamos el ángulo desde valores pequeños:

\begin{itemize}
    \item disminuye la componente horizontal;
    \item aumenta inicialmente la componente vertical;
    \item aumenta la altura máxima;
    \item aumenta el tiempo de vuelo, en el caso simétrico;
    \item el alcance aumenta hasta \(45\grado\) y luego disminuye, bajo las
    condiciones ideales del mismo nivel.
\end{itemize}

\begin{center}
\begin{tikzpicture}[scale=0.88]

    \draw[-{Latex},thick] (-0.3,0) -- (10.0,0)
        node[right] {\(x\)};

    \draw[-{Latex},thick] (0,-0.3) -- (0,5.0)
        node[above] {\(y\)};

    \draw[
        very thick,
        AzulSecundario,
        domain=0:8.7,
        samples=80
    ]
    plot (\x,{0.35*\x-0.04*\x*\x});

    \draw[
        very thick,
        VerdePrincipal,
        domain=0:9.2,
        samples=80
    ]
    plot (\x,{0.95*\x-0.103*\x*\x});

    \draw[
        very thick,
        RojoPrincipal,
        domain=0:8.7,
        samples=80
    ]
    plot (\x,{1.65*\x-0.19*\x*\x});

    \node[AzulSecundario] at (5.2,0.55) {ángulo pequeño};
    \node[VerdePrincipal] at (5.4,2.0) {ángulo intermedio};
    \node[RojoPrincipal] at (3.6,4.0) {ángulo grande};

\end{tikzpicture}
\end{center}

El dibujo es cualitativo: su función es mostrar que una misma rapidez inicial
puede producir trayectorias diferentes dependiendo del ángulo.


\section{Interpretación mediante vectores de velocidad}

Durante el vuelo:

\[
\vect{v}
=
v_x\ihat+v_y\jhat.
\]

La componente horizontal:

\[
v_x
\]

permanece constante.

En cambio:

\[
v_y
\]

disminuye linealmente con el tiempo.

Podemos pensar en tres momentos.

\subsection{Durante el ascenso}

\[
v_x>0,
\qquad
v_y>0.
\]

La velocidad apunta hacia adelante y hacia arriba.


\subsection{En el punto más alto}

\[
v_x>0,
\qquad
v_y=0.
\]

La velocidad es horizontal.


\subsection{Durante el descenso}

\[
v_x>0,
\qquad
v_y<0.
\]

La velocidad apunta hacia adelante y hacia abajo.

En los tres casos:

\[
\vect{a}=-g\jhat.
\]


\section{¿Por qué el proyectil no cae verticalmente de inmediato?}

Una confusión frecuente consiste en imaginar que, como la aceleración apunta
hacia abajo, el proyectil debería moverse únicamente hacia abajo.

Esto no ocurre porque aceleración y velocidad son magnitudes diferentes.

El proyectil posee una componente horizontal de velocidad:

\[
v_x=v_{0x}.
\]

Como:

\[
a_x=0,
\]

esa componente no desaparece.

Simultáneamente, la gravedad cambia:

\[
v_y.
\]

La combinación de ambos movimientos produce la trayectoria curva.

\begin{idea}
La aceleración indica cómo cambia la velocidad, no necesariamente la dirección
instantánea del movimiento.

Un proyectil puede tener:

\[
\vect{a}
\]

vertical hacia abajo y, al mismo tiempo:

\[
\vect{v}
\]

oblicua.
\end{idea}


\section{Estrategia general para resolver problemas}

\begin{estrategia}[Movimiento de proyectiles]
Para resolver un problema de proyectiles conviene seguir esta secuencia:

\begin{enumerate}
    \item Identifique el intervalo durante el cual el objeto puede modelarse como
    proyectil.

    \item Dibuje un esquema de la trayectoria.

    \item Defina claramente:
    \[
    x,\qquad y,\qquad O,\qquad +x,\qquad +y.
    \]

    \item Registre:
    \[
    x_0,\qquad y_0,\qquad x_f,\qquad y_f.
    \]

    \item Descomponga la velocidad inicial:
    \[
    v_{0x}=v_0\cos\theta_0,
    \]
    \[
    v_{0y}=v_0\sen\theta_0.
    \]

    \item Escriba las aceleraciones:
    \[
    a_x=0,
    \qquad
    a_y=-g.
    \]

    \item Analice el eje horizontal:
    \[
    x=x_0+v_{0x}t.
    \]

    \item Analice el eje vertical:
    \[
    y=y_0+v_{0y}t-\frac{1}{2}gt^2.
    \]

    \item Utilice el mismo tiempo en ambos ejes.

    \item Determine primero la variable que conecta los dos movimientos, que con
    frecuencia es \(t\).

    \item Reconstruya la velocidad final mediante:
    \[
    \vect{v}=v_x\ihat+v_y\jhat.
    \]

    \item Si se pide rapidez:
    \[
    v=\sqrt{v_x^2+v_y^2}.
    \]

    \item Si se pide dirección, determine el ángulo a partir de las componentes y
    compruebe el cuadrante.

    \item Verifique que cualquier fórmula especial utilizada cumpla sus
    condiciones de aplicación.

    \item Interprete físicamente el resultado.
\end{enumerate}
\end{estrategia}


\section{Cómo elegir la ecuación adecuada}

No es necesario utilizar todas las ecuaciones en cada problema.

\subsection{Cuando se busca el tiempo de vuelo}

La ecuación vertical suele ser la más útil:

\[
y_f
=
y_0
+
v_{0y}t
-
\frac{1}{2}gt^2.
\]


\subsection{Cuando se busca la posición horizontal}

Una vez conocido el tiempo:

\[
x_f
=
x_0+v_{0x}t.
\]


\subsection{Cuando se busca la altura máxima}

Utilice:

\[
v_y=0
\]

en el punto máximo.

Puede resultar eficiente usar:

\[
v_y^2
=
v_{0y}^2
-
2g(y-y_0).
\]


\subsection{Cuando se busca la velocidad de impacto}

Determine por separado:

\[
v_x
\]

y:

\[
v_y.
\]

Después reconstruya:

\[
\vect{v}
=
v_x\ihat+v_y\jhat.
\]


\section{Casos especiales}

Las diferentes situaciones habituales no requieren teorías distintas.

Todas proceden de:

\[
a_x=0,
\qquad
a_y=-g.
\]


\subsection{Lanzamiento horizontal}

\[
v_{0y}=0.
\]


\subsection{Lanzamiento oblicuo hacia arriba}

\[
v_{0x}>0,
\qquad
v_{0y}>0.
\]


\subsection{Lanzamiento oblicuo hacia abajo}

Con \(+y\) hacia arriba:

\[
v_{0x}>0,
\qquad
v_{0y}<0.
\]


\subsection{Lanzamiento vertical}

Si:

\[
v_{0x}=0,
\]

el movimiento deja de ser bidimensional y se reduce al movimiento vertical
estudiado anteriormente.

\begin{interpretacion}
Los casos especiales deben entenderse como consecuencias del modelo general.

No es necesario memorizar un conjunto diferente de ecuaciones para cada tipo de
lanzamiento.
\end{interpretacion}


\section{Casos límite}

Los casos límite permiten comprobar la coherencia de las ecuaciones.


\subsection{Ángulo nulo}

Si:

\[
\theta_0=0,
\]

entonces:

\[
v_{0x}=v_0,
\qquad
v_{0y}=0.
\]

Recuperamos el lanzamiento horizontal.


\subsection{Ángulo de noventa grados}

Si:

\[
\theta_0=90\grado,
\]

entonces:

\[
v_{0x}=0,
\qquad
v_{0y}=v_0.
\]

El movimiento se reduce a un lanzamiento vertical.


\subsection{Gravedad nula}

Si hipotéticamente:

\[
g\to0,
\]

entonces:

\[
a_x=a_y=0.
\]

El proyectil se movería con velocidad vectorial constante y su trayectoria sería
rectilínea:

\[
\vect{r}
=
\vect{r}_0+\vect{v}_0t.
\]


\section{Limitaciones del modelo ideal}

El modelo parabólico es muy útil, pero no describe perfectamente todos los
proyectiles reales.

\subsection{Resistencia del aire}

La resistencia del aire puede:

\begin{itemize}
    \item reducir la rapidez;
    \item disminuir la altura máxima;
    \item reducir el alcance;
    \item modificar la simetría de la trayectoria;
    \item hacer que la trayectoria deje de ser parabólica.
\end{itemize}


\subsection{Viento}

El movimiento del aire puede introducir modificaciones horizontales y
verticales que el modelo ideal no contempla.


\subsection{Rotación del proyectil}

Una pelota en rotación puede experimentar desviaciones adicionales asociadas a
su interacción con el aire.

Por ello, trayectorias de fútbol, béisbol, tenis o golf pueden diferir
considerablemente de una parábola ideal.


\subsection{Grandes distancias}

Para proyectiles de gran alcance pueden dejar de ser despreciables:

\begin{itemize}
    \item la curvatura terrestre;
    \item la rotación terrestre;
    \item la variación de \(g\) con la posición.
\end{itemize}

\begin{interpretacion}
Un modelo ideal no es ``incorrecto'' porque ignore ciertos efectos.

Es útil cuando conserva los factores dominantes para la precisión requerida.

La pregunta física correcta es:

\begin{quote}
¿son razonables las aproximaciones para esta situación?
\end{quote}
\end{interpretacion}


\section{Aplicaciones}

El movimiento de proyectiles aparece en una gran variedad de contextos.

\begin{itemize}
    \item deportes;
    \item chorros de agua;
    \item transporte de materiales;
    \item sistemas de descarga;
    \item experimentos de laboratorio;
    \item diseño de trayectorias idealizadas;
    \item análisis de objetos liberados desde vehículos.
\end{itemize}

\begin{aplicacion}[Chorro de agua]
Si un chorro sale de una boquilla con una velocidad inicial bien definida y los
efectos del aire son pequeños, cada pequeña porción de agua puede aproximarse
como un proyectil después de abandonar la boquilla.

La forma parabólica visible del chorro constituye una representación directa de
la trayectoria prevista por el modelo ideal.
\end{aplicacion}


\section{Comprobación experimental}

El modelo puede estudiarse mediante una cámara que registre posiciones en
intervalos iguales de tiempo.

Si se obtienen coordenadas:

\[
(x_i,y_i)
\]

para tiempos igualmente espaciados, el modelo ideal predice que:

\begin{itemize}
    \item los incrementos horizontales son aproximadamente iguales;
    \item la velocidad horizontal es aproximadamente constante;
    \item la coordenada vertical presenta comportamiento cuadrático con el tiempo;
    \item la velocidad vertical cambia linealmente con el tiempo;
    \item la aceleración vertical es aproximadamente constante.
\end{itemize}

Una gráfica:

\[
x-t
\]

debe ser aproximadamente lineal.

Una gráfica:

\[
y-t
\]

debe ser parabólica.

Una gráfica:

\[
v_y-t
\]

debe ser lineal con pendiente:

\[
-g.
\]

\begin{interpretacion}
Las desviaciones sistemáticas respecto de estas predicciones pueden indicar que
la resistencia del aire u otros efectos no son despreciables.
\end{interpretacion}


\section{Errores frecuentes}

\begin{errorfrecuente}
\textbf{Utilizar una sola ecuación escalar para todo el movimiento.}

El movimiento es bidimensional.

Debe analizarse mediante componentes:

\[
x
\qquad\text{y}\qquad
y.
\]
\end{errorfrecuente}


\begin{errorfrecuente}
\textbf{Asignar aceleración horizontal debido a que el proyectil avanza.}

Que exista velocidad horizontal no implica aceleración horizontal.

Dentro del modelo ideal:

\[
a_x=0.
\]
\end{errorfrecuente}


\begin{errorfrecuente}
\textbf{Creer que la velocidad horizontal disminuye durante el vuelo.}

Sin resistencia del aire:

\[
v_x=v_{0x}.
\]

La componente horizontal permanece constante.
\end{errorfrecuente}


\begin{errorfrecuente}
\textbf{Creer que en el punto más alto la velocidad total es cero.}

En general:

\[
v_y=0,
\]

pero:

\[
v_x\neq0.
\]

Por tanto:

\[
\vect{v}\neq\vect{0}.
\]
\end{errorfrecuente}


\begin{errorfrecuente}
\textbf{Creer que la aceleración es cero en el punto más alto.}

Durante todo el vuelo ideal:

\[
\vect{a}=-g\jhat.
\]
\end{errorfrecuente}


\begin{errorfrecuente}
\textbf{Usar tiempos diferentes para el movimiento horizontal y vertical.}

Existe un único movimiento físico y un único tiempo:

\[
t.
\]

Las componentes se evalúan simultáneamente.
\end{errorfrecuente}


\begin{errorfrecuente}
\textbf{Utilizar el alcance especial cuando las alturas son diferentes.}

La expresión:

\[
R
=
\frac{
v_0^2\sen(2\theta_0)
}{
g
}
\]

requiere que el proyectil regrese a la altura de lanzamiento.
\end{errorfrecuente}


\begin{errorfrecuente}
\textbf{Afirmar que \(45\grado\) siempre produce el alcance máximo.}

El resultado depende de las hipótesis:

\begin{itemize}
    \item ausencia de resistencia del aire;
    \item misma altura inicial y final;
    \item aceleración gravitatoria uniforme.
\end{itemize}
\end{errorfrecuente}


\begin{errorfrecuente}
\textbf{Confundir trayectoria y gráfica posición--tiempo.}

La trayectoria parabólica es una gráfica espacial:

\[
y\text{ contra }x.
\]

No es una gráfica:

\[
y\text{ contra }t.
\]

La función \(y(t)\) también es parabólica en este caso, pero representa una
relación física diferente.
\end{errorfrecuente}


\begin{errorfrecuente}
\textbf{Introducir el ángulo directamente en todas las ecuaciones sin
descomponer la velocidad.}

El procedimiento más seguro es calcular primero:

\[
v_{0x}=v_0\cos\theta_0,
\]

\[
v_{0y}=v_0\sen\theta_0.
\]

Después se trabaja por componentes.
\end{errorfrecuente}


\begin{errorfrecuente}
\textbf{Ignorar el signo de la componente vertical final.}

Durante el descenso:

\[
v_y<0
\]

si hemos elegido \(+y\) hacia arriba.

Una rapidez positiva no sustituye la información vectorial de la velocidad.
\end{errorfrecuente}


\section{Secuencia sugerida para impartir la clase}

\begin{notapedagogica}
Una posible secuencia para desarrollar este tema en clase es:

\begin{enumerate}
    \item Lanzar o imaginar una pelota de manera oblicua y preguntar qué ocurre
    simultáneamente en las direcciones horizontal y vertical.

    \item Recordar brevemente el movimiento en dos dimensiones y el método de
    componentes rectangulares.

    \item Definir qué intervalo del movimiento se modelará como movimiento de
    proyectil.

    \item Establecer explícitamente las hipótesis:
    \begin{itemize}
        \item resistencia del aire despreciable;
        \item \(g\) constante;
        \item superficie terrestre localmente plana.
    \end{itemize}

    \item Dibujar el sistema de referencia con:
    \[
    +x
    \]
    horizontal y:
    \[
    +y
    \]
    hacia arriba.

    \item Establecer:
    \[
    a_x=0,
    \qquad
    a_y=-g.
    \]

    \item Descomponer la velocidad inicial:
    \[
    v_{0x}=v_0\cos\theta_0,
    \]
    \[
    v_{0y}=v_0\sen\theta_0.
    \]

    \item Analizar primero el movimiento horizontal como MRU.

    \item Analizar después el movimiento vertical como caída libre.

    \item Insistir en que ambos movimientos comparten el mismo tiempo.

    \item Construir las ecuaciones generales:
    \[
    x=x_0+v_{0x}t,
    \]
    \[
    y=y_0+v_{0y}t-\frac{1}{2}gt^2.
    \]

    \item Dibujar los vectores velocidad en varios puntos de la trayectoria.

    \item Preguntar antes de responder qué ocurre con:
    \[
    v_x,\qquad v_y,\qquad\vect{a}
    \]
    en la altura máxima.

    \item Derivar la ecuación parabólica de la trayectoria eliminando \(t\).

    \item Resolver un lanzamiento horizontal.

    \item Resolver un lanzamiento oblicuo que regrese al mismo nivel.

    \item Deducir altura máxima, tiempo de vuelo y alcance como casos particulares,
    no como fórmulas aisladas.

    \item Mostrar por qué \(45\grado\) maximiza el alcance únicamente bajo las
    condiciones del caso simétrico.

    \item Resolver un problema con alturas inicial y final diferentes para mostrar
    que el método general continúa funcionando.

    \item Finalizar comparando el modelo ideal con una trayectoria real afectada por
    resistencia del aire.
\end{enumerate}
\end{notapedagogica}


\section{Preguntas de comprensión}

\begin{enumerate}
    \item ¿Qué condiciones deben cumplirse para modelar un objeto como proyectil
    ideal?

    \item ¿En qué momento comienza el movimiento de proyectil de una pelota lanzada
    con la mano?

    \item ¿Por qué resulta conveniente utilizar componentes rectangulares?

    \item Si \(+y\) apunta hacia arriba, ¿cuáles son las componentes de la
    aceleración de un proyectil ideal?

    \item ¿Por qué la componente horizontal de la velocidad permanece constante?

    \item ¿Por qué la componente vertical de la velocidad cambia?

    \item ¿El movimiento horizontal afecta el tiempo de caída vertical dentro del
    modelo ideal?

    \item ¿Qué variable comparten el movimiento horizontal y el vertical?

    \item ¿Qué representa físicamente \(v_{0x}\)?

    \item ¿Qué representa físicamente \(v_{0y}\)?

    \item ¿Qué ocurre con \(v_y\) mientras el proyectil asciende?

    \item ¿Cuál es el valor de \(v_y\) en el punto más alto?

    \item ¿Cuál es el valor de \(a_y\) en el punto más alto?

    \item ¿La velocidad total es cero en la altura máxima?

    \item ¿En qué dirección apunta la velocidad en el ápice?

    \item ¿Por qué la trayectoria ideal es parabólica?

    \item ¿Qué diferencia existe entre una trayectoria \(y-x\) y una gráfica
    \(y-t\)?

    \item Si un objeto se lanza horizontalmente desde una plataforma, ¿su componente
    vertical inicial es cero?

    \item ¿Significa eso que su aceleración vertical también es cero?

    \item ¿De qué depende el tiempo de caída de un lanzamiento horizontal desde una
    altura fija?

    \item Si duplicamos la rapidez horizontal de un lanzamiento desde la misma
    altura, ¿qué ocurre con el tiempo de caída?

    \item ¿Qué ocurre con el alcance horizontal en ese caso?

    \item ¿En qué condición el tiempo de descenso desde el punto más alto es igual
    al tiempo de ascenso?

    \item ¿Cuándo puede utilizarse:
    \[
    t_{\mathrm{vuelo}}
    =
    \frac{2v_0\sen\theta_0}{g}?
    \]

    \item ¿Cuándo puede utilizarse:
    \[
    R
    =
    \frac{v_0^2\sen(2\theta_0)}{g}?
    \]

    \item ¿Por qué los ángulos complementarios producen el mismo alcance bajo las
    condiciones ideales del caso simétrico?

    \item ¿Tienen también la misma altura máxima?

    \item ¿Por qué \(45\grado\) no debe considerarse universalmente el ángulo de
    alcance máximo?

    \item ¿Qué sucede con la trayectoria cuando la resistencia del aire deja de ser
    despreciable?

    \item ¿Puede un proyectil tener la misma rapidez en dos puntos diferentes de su
    trayectoria?

    \item Si el punto final está por debajo del punto inicial, ¿puede la rapidez de
    impacto ser mayor que la rapidez inicial?

    \item ¿Por qué una aceleración vertical puede producir una trayectoria curva
    aunque la componente horizontal de la velocidad sea constante?
\end{enumerate}


\begin{resumen}
El movimiento de proyectiles es un movimiento bidimensional que, dentro del
modelo ideal, puede analizarse mediante componentes rectangulares.

Se supone que:

\begin{itemize}
    \item la resistencia del aire es despreciable;
    \item la aceleración gravitatoria es constante;
    \item el proyectil se representa como una partícula.
\end{itemize}

Con:

\[
+x
\]

horizontal y:

\[
+y
\]

hacia arriba:

\[
\boxed{
a_x=0,
\qquad
a_y=-g
}.
\]

Si la rapidez inicial es \(v_0\) y el ángulo de lanzamiento es \(\theta_0\):

\[
\boxed{
v_{0x}=v_0\cos\theta_0
},
\]

\[
\boxed{
v_{0y}=v_0\sen\theta_0
}.
\]

El movimiento horizontal es uniforme:

\[
\boxed{
x=x_0+v_{0x}t
},
\]

\[
\boxed{
v_x=v_{0x}
}.
\]

El movimiento vertical es uniformemente acelerado:

\[
\boxed{
y=y_0+v_{0y}t-\frac{1}{2}gt^2
},
\]

\[
\boxed{
v_y=v_{0y}-gt
}.
\]

El vector velocidad es:

\[
\boxed{
\vect{v}
=
v_x\ihat+v_y\jhat
}.
\]

Su magnitud es:

\[
\boxed{
v=\sqrt{v_x^2+v_y^2}
}.
\]

El vector aceleración permanece:

\[
\boxed{
\vect{a}=-g\jhat
}.
\]

La ecuación de la trayectoria es:

\[
\boxed{
y
=
y_0
+
(x-x_0)\tan\theta_0
-
\frac{
g(x-x_0)^2
}{
2v_0^2\cos^2\theta_0
}
},
\]

por lo que la trayectoria ideal es parabólica.

En la altura máxima:

\[
\boxed{
v_y=0
},
\]

pero:

\[
\boxed{
v_x\neq0
}
\]

en general, y:

\[
\boxed{
a_y=-g
}.
\]

El tiempo hasta la altura máxima es:

\[
\boxed{
t_H
=
\frac{v_0\sen\theta_0}{g}
}.
\]

La altura ganada respecto del punto de lanzamiento es:

\[
\boxed{
H
=
\frac{
v_0^2\sen^2\theta_0
}{
2g
}
}.
\]

Si el proyectil regresa a la misma altura de lanzamiento:

\[
\boxed{
t_{\mathrm{vuelo}}
=
\frac{
2v_0\sen\theta_0
}{
g
}
},
\]

y:

\[
\boxed{
R
=
\frac{
v_0^2\sen(2\theta_0)
}{
g
}
}.
\]

Solamente bajo esas condiciones ideales, para una rapidez inicial fija, el
alcance máximo ocurre cuando:

\[
\boxed{
\theta_0=45\grado
}.
\]

Cuando las alturas inicial y final son diferentes, deben utilizarse las
ecuaciones generales por componentes.

La idea fundamental de todo el capítulo puede resumirse mediante:

\[
\boxed{
\begin{array}{c}
a_x=0
\\[2mm]
a_y=-g
\end{array}
}
\]

junto con:

\[
\boxed{
\text{un único tiempo para ambos movimientos}
}.
\]

Por ello, el procedimiento general consiste siempre en descomponer, analizar las
componentes rectangulares y reconstruir posteriormente las magnitudes vectoriales
que se necesiten.
\end{resumen}